% !TeX program = xelatex
% Shared simulation-specific extension; invoke from an analyses/* report directory.
% PGFPlots is loaded only for genuine scientific figures, never textual formulas.
% Shared XeLaTeX / KOMA-Script report style for composite_iq_checked_claims.
% Modeled on iq_test_differences_report.tex supplied by Anders Hellström.
% Requires XeLaTeX; loaded before \begin{document} from a report's main file.
% Reports keep their own mathematics, references, assumptions and AI appendix.
\documentclass[11pt,a4paper,oneside,open=any,parskip=half,headings=normal,
  numbers=noenddot,listof=totoc,bibliography=totoc]{scrbook}
\usepackage{iftex}
\RequireXeTeX
\usepackage[margin=18mm,top=19mm,bottom=20mm,headheight=16pt,headsep=7mm]{geometry}
\usepackage{fontspec}
\setmainfont{lmroman10-regular.otf}[BoldFont=lmroman10-bold.otf,
  ItalicFont=lmroman10-italic.otf,BoldItalicFont=lmroman10-bolditalic.otf]
\setsansfont{lmsans10-regular.otf}[BoldFont=lmsans10-bold.otf,
  ItalicFont=lmsans10-oblique.otf,BoldItalicFont=lmsans10-boldoblique.otf]
\setmonofont{DejaVu Sans Mono}[Scale=0.79]
\usepackage{mathtools,amsthm}
\usepackage{unicode-math}
\setmathfont{latinmodern-math.otf}
\usepackage[protrusion=true,expansion=false]{microtype}
\usepackage{xcolor}
\usepackage[most]{tcolorbox}
\usepackage{tabularray}
\UseTblrLibrary{booktabs}
\usepackage{booktabs,longtable,array}
\usepackage{graphicx,float,enumitem}
\usepackage{scrhack,listings}
\usepackage{needspace,placeins}
\usepackage[round,authoryear]{natbib}
\usepackage{xurl}
\usepackage{scrlayer-scrpage}
\usepackage{hyperref}
\usepackage{bookmark}
\usepackage[nameinlink,noabbrev]{cleveref}

% Centralized palette and link appearance from the reference report.
\definecolor{ink}{HTML}{183B56}
\definecolor{teal}{HTML}{007F82}
\definecolor{rust}{HTML}{B75D36}
\definecolor{pale}{HTML}{EFF5F7}
\newcommand{\ProjectAuthor}{Anders Hellström}
\hypersetup{unicode=true,colorlinks=true,linkcolor=ink,citecolor=teal,
  urlcolor=teal,pdfauthor={Anders Hellström; AI-assisted research and drafting}}
\clearpairofpagestyles
\automark[chapter]{chapter}
\renewcommand*{\chaptermarkformat}{}
\ihead{\small\sffamily \headmark}
\ohead{\small\sffamily \headmark}
\cfoot*{\pagemark}
\pagestyle{scrheadings}
\setkomafont{disposition}{\sffamily\bfseries\color{ink}}
\setkomafont{pageheadfoot}{\normalfont\sffamily\small\color{ink}}
\RedeclareSectionCommand[beforeskip=0pt,afterskip=14pt]{chapter}
\setlength{\emergencystretch}{2em}
\setlength{\parfillskip}{0pt plus 1fil}
\setcounter{tocdepth}{2}
\setcounter{secnumdepth}{2}
\setlist{nosep,leftmargin=*}
\setlength{\bibsep}{3pt}
\renewcommand{\bibfont}{\small}

% Preserve names and numbering used by the existing mathematical books.
\DeclareMathOperator{\E}{E}
\DeclareMathOperator{\Var}{Var}
\DeclareMathOperator{\Cov}{Cov}
\DeclareMathOperator{\SD}{SD}
\DeclareMathOperator{\Corr}{Corr}
\newcommand{\T}{\mathsf{T}}
\newcommand{\given}{\,\vert\,}
\newcommand{\one}{\mathbf{1}}
\newtheorem{theorem}{Theorem}[chapter]
\newtheorem{proposition}[theorem]{Proposition}
\newtheorem{lemma}[theorem]{Lemma}
\newtheorem{corollary}[theorem]{Corollary}
\crefname{chapter}{chapter}{chapters}
\Crefname{chapter}{Chapter}{Chapters}
\crefname{proposition}{proposition}{propositions}
\Crefname{proposition}{Proposition}{Propositions}
\crefname{corollary}{corollary}{corollaries}
\Crefname{corollary}{Corollary}{Corollaries}
\counterwithout{equation}{chapter}

% Boxes, code and front matter share the reference report's design.
\newtcolorbox{resultbox}[1][]{colback=pale,colframe=ink,boxrule=.7pt,
  arc=1mm,left=3mm,right=3mm,top=2mm,bottom=2mm,breakable,#1}
\newtcolorbox{disclosurebox}[1][]{colback=rust!7,colframe=rust,boxrule=1.1pt,
  arc=0mm,left=3mm,right=3mm,top=3mm,bottom=3mm,breakable,#1}
\lstset{basicstyle=\ttfamily\footnotesize,breaklines=true,columns=fullflexible,
  keepspaces=true,showstringspaces=false,numbers=left,numberstyle=\tiny\color{gray},
  numbersep=6pt,xleftmargin=15pt,frame=single,rulecolor=\color{ink!25},
  framerule=.3pt,framesep=5pt,backgroundcolor=\color{pale!45},
  keywordstyle=\color{ink}\bfseries,commentstyle=\color{teal},
  stringstyle=\color{rust},tabsize=2,upquote=true}
\newenvironment{abstract}{\addchap{Overview}}{}

% Report-specific summaries are supplied in each book's preamble.
\newcommand{\ProjectReportType}{MATHEMATICAL AND COMPUTATIONAL REPORT}
\newcommand{\ProjectReportHighlights}{\textbf{Scope:} mathematical deductions,
  numerical approximations and empirical claims are distinguished.}
\makeatletter
\newcommand*{\ProjectReportTitlePage}{%
  \begin{titlepage}
    \thispagestyle{empty}
    \vspace*{6mm}
    {\sffamily\color{teal}\large \ProjectReportType\par}
    \vspace{10mm}
    {\sffamily\bfseries\color{ink}\fontsize{27}{32}\selectfont\@title\par}
    \vspace{6mm}
    {\Large\@subtitle\par}
    \vspace{8mm}
    \begin{disclosurebox}[title={\sffamily\bfseries AI-GENERATED REPORT: READ THIS FIRST}]
      This document and substantial supporting mathematical exposition, programs,
      verification and typesetting were prepared with AI assistance at the request
      of Anders Hellström. The reported numerical results are conditional on
      explicit models and evidence; automated checks are not independent human
      peer review or empirical psychometric validation. Full AI-use details
      appear in the final appendix.
    \end{disclosurebox}
    \vspace{5mm}
    \begin{resultbox}[title={\sffamily\bfseries Key findings and interpretation}]
      \ProjectReportHighlights
    \end{resultbox}
    \vspace{3mm}
    \begin{resultbox}[title={\sffamily\bfseries Project repository and reproducibility}]
      \textbf{Repository:}\quad
      \url{https://github.com/AndersH3/composite_iq_checked_claims}

      Source, proof and simulation files are organized in the repository;
      this report specifies its own assumptions, references and verification scope.
    \end{resultbox}
    \vfill
    Prepared at the request of \textbf{\ProjectAuthor}\\
    Report date: \textbf{\@date}\\
    Authoring workflow: OpenAI-assisted research and typesetting\\
    Format: XeLaTeX, KOMA-Script \texttt{scrbook}; A4, narrow margins
  \end{titlepage}%
}
\makeatother

\usepackage{pgfplots}
\pgfplotsset{compat=1.18,width=.92\linewidth,height=6.0cm,
  every axis/.append style={grid=major,grid style={gray!20},
    tick label style={font=\small},label style={font=\small},
    legend style={font=\small,draw=none,fill=white}}}
\usepgfplotslibrary{groupplots}
\renewcommand{\ProjectReportType}{SIMULATION METHODS AND RESULTS}
\renewcommand{\ProjectReportRunningTitle}{Simulation methods and results}
\renewcommand{\ProjectReportRunningSubtitle}{Python/Dask experiments}

\title{Simulating Multiple-Test Admissions}
\subtitle{Python/Dask methods, results and verification limits}
\author{\ProjectAuthor}
\date{16 November 2026}

% Edition-specific findings displayed on the shared reference-style title page.
\renewcommand{\ProjectReportHighlights}{Seeded Python/Dask experiments estimate union, intersection and
conditional event frequencies for correlated Gaussian tests and compare
them with reference integrals. The reported counts and uncertainty
do not establish actual membership frequencies.}

\begin{document}
\frontmatter
\ProjectReportTitlePage
% Computed by the Dask report task; display rounding only.
\newcommand{\SimTrials}{10\,000}
\newcommand{\SimReplicates}{1}
\newcommand{\SimPooledTrials}{10\,000}
\newcommand{\SimChunk}{2\,500}
\newcommand{\SimMaxMeanSE}{3.166}
\newcommand{\SimMaxVarianceRatio}{0.000000}
\newcommand{\SimRareMin}{0}
\newcommand{\SimRareMax}{0}
\newcommand{\SimMaxEventSE}{3.166}
\newcommand{\SimMaxSDRelativeError}{0.000000}


\begin{abstract}
This report describes the Python/Dask simulation of multiple-test admission
rules under a specified correlated Gaussian model. A Dask task graph
generates shared correlated pairs, evaluates four threshold rules and a
fifth calibrated union rule, reduces independently seeded replicates,
and compares event, union and conditional frequencies with numerical
Gaussian reference integrals. We state denominators, rare-event
uncertainty and finite-model limitations explicitly. Deal contracts
and runtime tests address selected implementation properties; the
simulation does not prove independence of pseudorandom observations
or establish membership frequencies in a real population.
\end{abstract}

\noindent\textbf{AI disclosure.} OpenAI Codex and collaborating AI agents
prepared the exposition, Dask analysis, source code and typesetting within the
repository's disclosed workflow. Independent human peer review is not claimed.
The numerical tables are generated by executed code rather than invented for
the narrative. Full disclosure appears in \cref{sec:ai-disclosure}.

\noindent\textbf{Proof coverage.} Appendices~\ref{sec:proof-foundations},
\ref{sec:shared-statistical-methods} and \ref{sec:towers-proofs} provide full
derivations from referenced theorems. The complete claim register in
Appendix~\ref{sec:towers-simulation-claim-map} distinguishes mathematical
deductions, rounded numerical evaluations and retained execution records.

\tableofcontents
\mainmatter
\chapter{The experiment and the questions it answers}

The accompanying main analysis jointly assesses Towers's original letter and
its follow-up \citep{towersOriginal,towersFollowup}. The follow-up retracts the
original overlap calculation. This companion report concentrates on the
simulation experiment and its interpretation; a simulation does not convert
an incorrect probability identity into a correct one, nor does numerical
agreement establish that the Gaussian model describes actual examinees.

Each simulated observation is a pair of standardized scores \((X,Y)\), with
population means zero, population variances one, and correlation \(0.7\).
For cutoffs \(a,b\), write
\[
 A=\{X>a\},\qquad B=\{Y>b\},\qquad
 p_A=P(A),\quad p_B=P(B),\quad j=P(A\cap B).
\]
The main output quantities are
\begin{align}
 u&=P(A\cup B)=p_A+p_B-j,\label{eq:towers-union}\\
 q_{B\mid A}&=P(B\mid A)=\frac{j}{p_A},\qquad
 q_{A\mid B}=P(A\mid B)=\frac{j}{p_B}.
 \label{eq:towers-conditionals}
\end{align}
The union concerns a rule accepting either score. A conditional probability
concerns passers of one specified test. The joint probability concerns passing
both tests in the whole reference population. Reporting one as another would
answer a different question even if every simulated count were exact.
Proposition~\ref{tw:events} derives each event formula and denominator.


\noindent\textbf{Block-aggregated Dask experiment.}
The source supports a Dask Array reference engine and an optimized
Dask delayed engine. The latter evaluates all five decision rules
on a shared Gaussian pair block and returns only integer four-cell
counts to the scheduler. The numerical records and provenance manifest,
rather than the source defaults, determine the number of executed
observations and the actual engine for this edition.

The study uses \SimReplicates{} replicate streams, each with
\SimTrials{} pairs, for \SimPooledTrials{} pairs per rule after pooling.
Dask processes blocks of \SimChunk{} observations. The same pairs are
classified under all five rules within a replicate. This common-random-number
design permits direct rule comparisons and avoids generating the same
population separately for every cutoff. It also makes the rule-specific
estimates statistically dependent; the five rows are not five independent
replications of the same test. Separate seeded replicate streams provide the
between-run comparison. Their pseudorandom construction and the finite
sampling assumptions remain distinct from a mathematical independence proof.

\section{The five decision rules}

Four rules reproduce the questions in the existing analysis. The last asks a
new implementation question using its already calculated cutoff.
\begin{center}
\begin{tabular}{llll}
\toprule
Label & Rule & Marginal tail on \(X\) & Marginal tail on \(Y\)\\
\midrule
E2 & Equal top 2\% & \(0.02\) & \(0.02\)\\
E0.1 & Equal top 0.1\% & \(0.001\) & \(0.001\)\\
Rare & Equal top 1 in 33,000 & \(1/33000\) & \(1/33000\)\\
Mixed & Top 2\% and top 0.1\% & \(0.02\) & \(0.001\)\\
OR2 & Calibrated equal-cutoff OR & determined by \(u=0.02\) & same\\
\bottomrule
\end{tabular}
\end{center}
For a marginal tail probability \(p\), the cutoff is
\(\Phi^{-1}(1-p)\), calculated using a survival-function inverse to avoid
forming a small tail by subtracting two nearly equal floating-point numbers.
Strict and non-strict cutoffs have the same theoretical probability for the
continuous Gaussian distribution. The executable nevertheless uses one
consistent strict comparison for every event count.

\chapter{How the simulated pairs are constructed}

For each observation the seeded Gaussian model supplies two nominally independent
standard normal values \(G,E\). The transformed pair is
\begin{equation}
 X=G,\qquad Y=\rho G+\sqrt{1-\rho^2}\,E,\qquad \rho=0.7.
 \label{eq:towers-generator}
\end{equation}
This construction replaces a general covariance factorization by a transparent
two-dimensional formula. Theoretical mean, variance and covariance checks are
immediate:
\[
 E[X]=E[Y]=0,\qquad \operatorname{Var}(X)=1,
 \qquad \operatorname{Var}(Y)=\rho^2+(1-\rho^2)=1,
 \qquad \operatorname{Cov}(X,Y)=\rho.
\]
These equalities follow by linearity, covariance bilinearity and the assumed
independence of \(G,E\). An affine transformation of a joint Gaussian vector
is Gaussian \citep[Chapter 3]{gallager2013}; therefore the construction has
the required bivariate normal law in the ideal independent-normal model.
It does not prove that a finite deterministic generator has that exact law.
Proposition~\ref{tw:gaussian} supplies the density change of variables and
the complete conditional-law derivation under the ideal assumptions.

The simulation classifies each pair into exactly one of four cells:
\[
 \neg A\cap\neg B,\quad A\cap\neg B,\quad
 \neg A\cap B,\quad A\cap B.
\]
The corresponding counts are denoted
\(n_{00},n_{10},n_{01},n_{11}\), with total \(T\). Their exact integer
identities are
\begin{align*}
 n_{00}+n_{10}+n_{01}+n_{11}&=T,\\
 n_A=n_{10}+n_{11},\qquad n_B&=n_{01}+n_{11},\\
 n_U=n_{10}+n_{01}+n_{11}&=n_A+n_B-n_{11}.
\end{align*}
These hold because the four cells partition the sample. They are checked as
bookkeeping identities rather than assessed with a statistical tolerance.
The distinction matters: random variation can explain an empirical joint
frequency differing from its theoretical target, but cannot explain cell
counts failing to sum to the number of observations.
The pointwise proof is Proposition~\ref{tw:simulation}.

\section{The numerical targets are calculated separately}

The reference joint tail is the one-dimensional integral
\begin{equation}
 j(a,b;\rho)=\int_a^\infty
 \phi(x)\,\overline\Phi\!\left(
 \frac{b-\rho x}{\sqrt{1-\rho^2}}\right)\,dx.
 \label{eq:towers-reference-integral}
\end{equation}
Indeed, \eqref{eq:towers-generator} makes \(Y\mid X=x\) normal with mean
\(\rho x\) and variance \(1-\rho^2\). Integrating its exceedance probability
against the standard normal density of \(X\) proves the formula by the law of
total probability. The repository's existing numerical function also compares
this integral with integration of the bivariate density over the correlation
parameter. The Dask graph schedules that maintained SciPy calculation as a
separate task. This is an independent algorithmic comparison with event
counting; it is not a second simulation of the same random pairs.

Inclusion-exclusion proves \eqref{eq:towers-union} by counting the intersection
once rather than twice. The conditional-probability definition on events of
positive probability proves \eqref{eq:towers-conditionals}. Thus the target
relationships do not depend on a Monte Carlo result or on a graphical reading
of overlapping normal curves.
Propositions~\ref{tw:gaussian} and \ref{tw:plackett} prove both integration
identities, including the correlation-derivative identity and the required
justification for differentiating under the integral.

\chapter{The pooled results}
\label{sec:towers-pooled}

\section{Joint events and eligibility through either test}

The first table reports the integer evidence and the joint and union
frequencies. Probabilities are fractions of the entire simulated population.
The joint and union columns therefore have the same total denominator,
\(\SimPooledTrials\), while describing different events.

\begingroup\small
\setlength{\tabcolsep}{4pt}
\Needspace{6\baselineskip}
\noindent\textbf{Pooled event-cell counts}
\begin{longtable}{rrrrr}
\toprule
Rule & Neither & $A$ only & $B$ only & Both \\
\midrule
\endfirsthead
Rule & Neither & $A$ only & $B$ only & Both \\
\midrule
\endhead
\bottomrule
\endfoot
E2 & $9653$ & $146$ & $150$ & $51$ \\
E0.1 & $9982$ & $4$ & $12$ & $2$ \\
Rare & $9997$ & $1$ & $2$ & $0$ \\
Mixed & $9795$ & $191$ & $8$ & $6$ \\
Calibrated & $9795$ & $85$ & $96$ & $24$ \\
\end{longtable}
\endgroup

\begingroup\small
\setlength{\tabcolsep}{4pt}
\Needspace{6\baselineskip}
\noindent\textbf{Joint probabilities and 95\% Wilson bounds}
\begin{longtable}{rrrrr}
\toprule
Rule & Target & Observed & 95\% lower & 95\% upper \\
\midrule
\endfirsthead
Rule & Target & Observed & 95\% lower & 95\% upper \\
\midrule
\endhead
\bottomrule
\endfoot
E2 & $0.006276$ & $0.006247$ & $0.006170$ & $0.006325$ \\
E0.1 & $1.595892\times10^{-4}$ & $1.570000\times10^{-4}$ & $1.451925\times10^{-4}$ & $1.697676\times10^{-4}$ \\
Rare & $2.312498\times10^{-6}$ & $2.750000\times10^{-6}$ & $1.535610\times10^{-6}$ & $4.924749\times10^{-6}$ \\
Mixed & $6.585888\times10^{-4}$ & $6.567500\times10^{-4}$ & $6.321191\times10^{-4}$ & $6.823400\times10^{-4}$ \\
OR2 & $0.003204$ & $0.003181$ & $0.003126$ & $0.003236$ \\
\end{longtable}
\endgroup


\begingroup\small
\setlength{\tabcolsep}{4pt}
\Needspace{6\baselineskip}
\noindent\textbf{OR-rule probabilities and 95\% Wilson bounds}
\begin{longtable}{rrrrr}
\toprule
Rule & Target & Observed & 95\% lower & 95\% upper \\
\midrule
\endfirsthead
Rule & Target & Observed & 95\% lower & 95\% upper \\
\midrule
\endhead
\bottomrule
\endfoot
E2 & $0.033724$ & $0.033634$ & $0.033458$ & $0.033811$ \\
E0.1 & $0.001840$ & $0.001801$ & $0.001760$ & $0.001843$ \\
Rare & $5.829356\times10^{-5}$ & $6.150000\times10^{-5}$ & $5.428017\times10^{-5}$ & $6.968007\times10^{-5}$ \\
Mixed & $0.020341$ & $0.020343$ & $0.020205$ & $0.020482$ \\
OR2 & $0.020000$ & $0.019815$ & $0.019679$ & $0.019952$ \\
\end{longtable}
\endgroup


\begingroup\small
\setlength{\tabcolsep}{4pt}
\Needspace{6\baselineskip}
\noindent\textbf{Rare-rule joint probability: exact binomial interval}
\begin{longtable}{rrrr}
\toprule
Joint target & Observed & Exact 95\% lower & Exact 95\% upper \\
\midrule
\endfirsthead
Joint target & Observed & Exact 95\% lower & Exact 95\% upper \\
\midrule
\endhead
\bottomrule
\endfoot
$2.312498\times10^{-6}$ & $0$ & $0$ & $3.688199\times10^{-4}$ \\
\end{longtable}
\endgroup



The equal-top-2\% case is the most informative comparison with the familiar
``2\%'' description. The Gaussian reference predicts
\(j\approx0.0062762740\) and
\(u\approx0.0337237260\): approximately 0.628\% pass both tests and
3.372\% pass at least one. The latter is larger than 2\% because a positive
fraction qualify on just one test. Positive correlation creates substantial
overlap, but does not make the two passing events identical.

For the equal-top-0.1\% case the targets are
\(j\approx0.0001595892\) and
\(u\approx0.0018404108\). The extra route to qualification again enlarges
the eligible population, although the absolute probabilities are smaller.
The sample now contains fewer joint successes, so relative uncertainty is
larger than for E2 even though the total number of simulated pairs is the same.
The rare case magnifies this difference between absolute accuracy and
relative precision; it is examined separately below.

The mixed rule is asymmetric. Its reference union is
\(u\approx0.0203414112\), only modestly above 2\%, because most passers of
the much higher cutoff are already included among passers of the lower one.
That observation is about the modeled events, not a claim that one actual
test subsumes another or that their errors behave in this way empirically.

% Plot coordinates are prepared in Dask; PGFPlots only renders the CSV columns.
\begin{figure}[htbp]
\centering
\begin{tikzpicture}
\begin{axis}[
 width=0.48\textwidth,height=5.6cm,
 ymode=log,ylabel={Joint probability},xlabel={Rule},
 xtick={1,2,3,4,5},xticklabels={E2,E0.1,Rare,Mixed,OR2},
 xticklabel style={font=\small},grid=both,
 legend style={font=\small,at={(0.5,1.02)},anchor=south,legend columns=2},
 xmin=0.6,xmax=5.4]
\addplot+[only marks,mark=*,blue] table[x=case_index,y=joint,col sep=comma]{simulation-data/events.csv};
\addlegendentry{Simulation}
\addplot+[only marks,mark=x,black] table[x=case_index,y=theory_joint,col sep=comma]{simulation-data/events.csv};
\addlegendentry{Gaussian target}
\end{axis}
\end{tikzpicture}\hfill
\begin{tikzpicture}
\begin{axis}[
 width=0.48\textwidth,height=5.6cm,
 ylabel={OR eligibility probability},xlabel={Rule},
 xtick={1,2,3,4,5},xticklabels={E2,E0.1,Rare,Mixed,OR2},
 xticklabel style={font=\small},grid=major,
 xmin=0.6,xmax=5.4,ymin=0]
\addplot+[only marks,mark=*,blue] table[x=case_index,y=union,col sep=comma]{simulation-data/events.csv};
\addplot+[only marks,mark=x,black] table[x=case_index,y=theory_union,col sep=comma]{simulation-data/events.csv};
\end{axis}
\end{tikzpicture}
\caption{Empirical probabilities and separately computed Gaussian targets.
The joint-probability axis is logarithmic because the cases span several
orders of magnitude. Confidence intervals and actual event counts are in the
tables; close plotting positions alone do not establish precision.}
\label{fig:towers-events}
\end{figure}

\FloatBarrier
\section{Repeat qualification and its denominator}

For a positive simulated count \(n_A\), the empirical repeat-qualification
rate is \(n_{11}/n_A\), not \(n_{11}/T\). Likewise the reverse conditional
uses \(n_B\). The conditional table retains the actual denominators, so a
rate based on a small subgroup cannot be mistaken for one based on all
\(\SimPooledTrials\) observations.

\begingroup\small
\setlength{\tabcolsep}{4pt}
\Needspace{6\baselineskip}
\noindent\textbf{Repeat qualification among passers of test A}
\begin{longtable}{rrrrrr}
\toprule
Rule & $N_A$ & $P(B|A)$ target & Observed & 95\% lower & 95\% upper \\
\midrule
\endfirsthead
Rule & $N_A$ & $P(B|A)$ target & Observed & 95\% lower & 95\% upper \\
\midrule
\endhead
\bottomrule
\endfoot
E2 & $80075$ & $0.313814$ & $0.312057$ & $0.308857$ & $0.315276$ \\
E0.1 & $3908$ & $0.159589$ & $0.160696$ & $0.149516$ & $0.172543$ \\
Rare & $115$ & $0.076312$ & $0.095652$ & $0.054253$ & $0.163192$ \\
Mixed & $80075$ & $0.032929$ & $0.032807$ & $0.031595$ & $0.034063$ \\
OR2 & $46206$ & $0.276178$ & $0.275332$ & $0.271278$ & $0.279424$ \\
\end{longtable}
\endgroup

\begingroup\small
\setlength{\tabcolsep}{4pt}
\Needspace{6\baselineskip}
\noindent\textbf{Repeat qualification among passers of test B}
\begin{longtable}{rrrrrr}
\toprule
Rule & $N_B$ & $P(A|B)$ target & Observed & 95\% lower & 95\% upper \\
\midrule
\endfirsthead
Rule & $N_B$ & $P(A|B)$ target & Observed & 95\% lower & 95\% upper \\
\midrule
\endhead
\bottomrule
\endfoot
E2 & $79449$ & $0.313814$ & $0.314516$ & $0.311297$ & $0.317754$ \\
E0.1 & $3925$ & $0.159589$ & $0.160000$ & $0.148864$ & $0.171801$ \\
Rare & $142$ & $0.076312$ & $0.077465$ & $0.043803$ & $0.133385$ \\
Mixed & $3925$ & $0.658589$ & $0.669299$ & $0.654422$ & $0.683846$ \\
OR2 & $45776$ & $0.276178$ & $0.277919$ & $0.273834$ & $0.282041$ \\
\end{longtable}
\endgroup



For equal thresholds, the two theoretical conditional probabilities agree by
symmetry, but their simulated values need not be identical: the observed
marginal counts \(n_A,n_B\) fluctuate separately. For E2, the target is
approximately \(31.3814\%\), rather than 2\% and rather than 0.628\%.
For E0.1 the repeat-qualification target is approximately \(15.9589\%\).
At the rare equal thresholds it is approximately \(7.6312\%\). These
differences illustrate that fixing the continuous correlation at 0.7 does
not fix repeat qualification at a single percentage independently of cutoff.

The asymmetry in the mixed case is especially useful. The target
\(P(B\mid A)\) is approximately \(3.29294\%\), while
\(P(A\mid B)\) is approximately \(65.8589\%\). Both use the same joint
probability, approximately \(0.0006585888\), divided by different marginal
probabilities, 0.02 and 0.001. The factor-of-20 difference follows directly
from those denominators. Reversing the conditioning event would change the
question rather than merely changing its wording.

For these subgroup rates, a confidence interval is conditional on the observed
number in the selected subgroup. Given \(n_A=m>0\), the both-pass count is
binomial with \(m\) trials and success probability \(q_{B\mid A}\) under
the independent-pair model. To see this, condition on the identities of the
\(m\) observations in \(A\): their \(B\) indicators are independent with
that conditional success probability. Their sum is therefore binomial.
Mixing over which identities were selected changes neither that conditional
law nor its denominator. If no observation passes \(A\), the conditional
rate is undefined; it must not be reported as zero.
Proposition~\ref{sm:conditional-count} gives a full conditional-binomial proof;
Proposition~\ref{tw:simulation} applies it to this four-cell sample.

\chapter{Rare events: what the counts can resolve}
\label{sec:towers-rare}

At the equal top-1-in-33,000 cutoffs, the model gives
\[
 j\approx2.3124979\times10^{-6},\qquad
 u\approx5.8293563\times10^{-5}.
\]
A replicate of \(\SimTrials\) pairs therefore has expected joint count
\(\SimTrials\times j\), while the expected number of passers of either
specified marginal test is \(\SimTrials/33{,}000\).
Pooling \(\SimReplicates\) such replicates increases the expected
joint count in direct proportion to the number of pairs; it does not
create millions of observations of the rare event itself.
The relevant information is the small joint count displayed in the table.

For a binomial count \(K\) from \(T\) independent pairs,
\(E[K]=Tj\) and \(\operatorname{Var}(K)=Tj(1-j)\), by the moment
identities for sums of independent indicators. Consequently the relative
standard error of the frequency is
\[
 \frac{\sqrt{j(1-j)/T}}{j}
 =\sqrt{\frac{1-j}{Tj}}.
\]
It is controlled mainly by the expected number of successes \(Tj\) when
\(j\) is small. A large total denominator can coexist with weak relative
precision for a rare numerator. The ideal probability of seeing zero joint
successes is \((1-j)^T\), obtained by multiplying the no-success probability
across independent pairs. A zero count therefore does not show that the
underlying joint probability is zero.
Proposition~\ref{sm:binomial} and Corollary~\ref{sm:rare} derive these moment, relative
error and zero-count expressions without a normal approximation to the count.

The pooled intervals and replicate counts quantify this limitation rather
than concealing it with extra decimal places. Wilson intervals summarize the
sampling uncertainty; the rare-event summary also provides an exact binomial
interval. ``Exact'' concerns inversion of the binomial distribution, not an
exact probability statement about real tests. Discreteness generally makes
its coverage conservative. Neither interval corrects a misspecified Gaussian
model, unrepresentative reference population, practice effect or dependence
between real attempts.

\section{Between-replicate results}

The following table preserves each replicate rather than only the pooled
average. Within a replicate the rule-specific counts share the same pairs;
between replicate streams, different seeds provide a reproducibility and
sampling-variation comparison.

\begingroup\small
\setlength{\tabcolsep}{4pt}
\Needspace{6\baselineskip}
\noindent\textbf{Each seeded replicate retains its own event denominators}
\begin{longtable}{rrrrrr}
\toprule
Replicate & Rule & $N_A$ & $N_B$ & Both & Union \\
\midrule
\endfirsthead
Replicate & Rule & $N_A$ & $N_B$ & Both & Union \\
\midrule
\endhead
\bottomrule
\endfoot
$1$ & E2 & $197$ & $201$ & $51$ & $347$ \\
$1$ & E0.1 & $6$ & $14$ & $2$ & $18$ \\
$1$ & Rare & $1$ & $2$ & $0$ & $3$ \\
$1$ & Mixed & $197$ & $14$ & $6$ & $205$ \\
$1$ & Calibrated & $109$ & $120$ & $24$ & $205$ \\
\end{longtable}
\endgroup


The rare-case joint count ranges from \SimRareMin{} to \SimRareMax{}
across the recorded replicates. This is variation among executions under the
same input model, not variation in the specified correlation. The largest
standardized event-frequency discrepancy in the recorded diagnostic set is
\SimMaxEventSE{} binomial standard errors. A diagnostic standardized by a
binomial standard error is informative about ordinary sampling variation,
but its normal approximation is coarse for expected counts near two. Exact
binomial calculations are more appropriate for interpreting that rare row.

Confidence intervals in this report are pointwise. A 95\% interval is not a
95\% simultaneous guarantee for every joint, union and conditional estimate
in every case and replicate. Common random numbers also mean that comparing
rows as if their errors were independent would give the wrong uncertainty
for a difference. The recorded broad computational acceptance checks are
separate from the displayed 95\% descriptive intervals; no seeds are selected
or replaced because an interval is inconvenient.
Propositions~\ref{sm:wilson}, \ref{sm:clopper-pearson}, \ref{sm:pooling} and
\ref{sm:multiple} prove the interval constructions, conservative exact
coverage, pooling and the limits of simultaneous interpretation.

\chapter{Calibrating a 2\% multiple-test rule}

E2 sets each individual passing probability to 2\%, but the combined OR rule
has target probability about 3.3724\%. OR2 instead asks for the common cutoff
\(t\) satisfying
\begin{equation}
 2\overline\Phi(t)-j(t,t;0.7)=0.02.
 \label{eq:towers-calibration}
\end{equation}
The previously calculated solution is approximately
\(t=2.2700549291\). On the main report's mean-100, SD-16 reporting scale,
this is approximately \(136.320879\), with a marginal tail probability of
about \(0.0116021252\) on each test. These are definitions within the
Gaussian example, not revised norms for actual instruments.

\begingroup\small
\setlength{\tabcolsep}{4pt}
\Needspace{6\baselineskip}
\noindent\textbf{Calibrated common cutoff and observed combined eligibility}
\begin{longtable}{rrrrrr}
\toprule
Cutoff $t$ & SD-16 score & Target OR & Observed OR & 95\% lower & 95\% upper \\
\midrule
\endfirsthead
Cutoff $t$ & SD-16 score & Target OR & Observed OR & 95\% lower & 95\% upper \\
\midrule
\endhead
\bottomrule
\endfoot
$2.270055$ & $136.320879$ & $0.02$ & $0.020500$ & $0.017901$ & $0.023467$ \\
\end{longtable}
\endgroup

% Promote the retained article heading while preserving the recorded data.
\begingroup
\let\subsection\section
\subsection{Specific observations in the recorded run}
The nominal E2 rule admits 347 of 10000 pairs ($3.4700$\%), while OR2 admits 205 ($2.0500$\%). The latter is $1.423$ binomial standard errors from its fixed model target. The target is inside this run's displayed 95\% Wilson interval. Such a pointwise interval can miss a known target in a correctly specified simulation; the much more conservative, predeclared implementation diagnostic still passes. The cutoff and seeds were not retuned.

For Rare, 0 joint passes occur among 1 passers of A. Thus the conditional estimate is based on that small actual denominator, not the full 10000 pairs. The conditional interval and the separate exact joint interval make the loss of rare-event precision visible.

\endgroup

There is a useful model-level existence and uniqueness argument. Increasing
\(t\) shrinks the event \(\{X>t\text{ or }Y>t\}\). Because the bivariate
normal density at \(|\rho|<1\) is strictly positive, moving from any finite
\(t_1\) to a larger finite \(t_2\) removes a region of positive probability;
the union probability therefore strictly decreases. The continuous density
gives continuity in \(t\), and the probability tends to one and zero at
negative and positive infinity, respectively. The intermediate value theorem
then gives a unique solution for every target strictly between zero and one.
Equation~\eqref{eq:towers-calibration} is the inclusion-exclusion expression
for that same probability. A numerical root finder approximates the unique
root; empirical counts do not define it.
Proposition~\ref{tw:calibration} provides the detailed proof, including
the positive-probability set difference and continuity argument.

The OR2 simulation holds that calculated cutoff fixed and estimates its
achieved union rate. It is a validation of the implemented decision rule under
the model, not a process of tuning the cutoff to make a particular random run
contain exactly 2\% passers. Conditional repeat qualification and joint passing
remain separate outputs even when the union is calibrated to 2\%.



\section{Continuous correlation is not binary correlation}

Passing indicators have phi correlation
\begin{equation}
 \phi_{A,B}=\frac{j-p_Ap_B}
 {\sqrt{p_A(1-p_A)p_B(1-p_B)}}.
 \label{eq:towers-phi}
\end{equation}
For indicators, \(E[1_A]=p_A\), \(\operatorname{Var}(1_A)=p_A(1-p_A)\)
and \(\operatorname{Cov}(1_A,1_B)=j-p_Ap_B\). Substituting those identities
into the Pearson-correlation definition proves the formula. It generally
differs from the continuous-score correlation \(\rho\); the E2 reference,
for example, is approximately 0.29981 although \(\rho=0.7\). The discrepancy
is a consequence of thresholding and the margins, not an error that simulation
should remove. Historical disagreement caused by confusing these correlations
cannot be repaired by changing the sample size.
Proposition~\ref{tw:events} proves the formula and its distinction from
continuous-score correlation.


\begingroup\small
\setlength{\tabcolsep}{4pt}
\Needspace{6\baselineskip}
\noindent\textbf{Continuous-score correlation and thresholded indicator correlation}
\begin{longtable}{rrrr}
\toprule
Rule & Continuous correlation & Binary phi target & Binary phi observed \\
\midrule
\endfirsthead
Rule & Continuous correlation & Binary phi target & Binary phi observed \\
\midrule
\endhead
\bottomrule
\endfoot
E2 & $0.7$ & $0.299810$ & $0.241195$ \\
E0.1 & $0.7$ & $0.158748$ & $0.217519$ \\
Rare & $0.7$ & $0.076284$ & $-1.414426\times10^{-4}$ \\
Mixed & $0.7$ & $0.144315$ & $0.110165$ \\
Calibrated & $0.7$ & $0.267681$ & $0.200711$ \\
\end{longtable}
\endgroup


\chapter{Dask execution, provenance and verification scope}

The new computation is expressed as a Dask graph: seeded array generation,
score transformation, event masks, integer reductions, replicate pooling,
reference calculations, interval evaluation and report-data serialization.
Dask schedules work by chunks, allowing the pair arrays to be reduced without
keeping all simulated pairs in one in-memory table. The pinned Dask/NumPy
versions, seeds, chunk sizes and graph execution settings are part of the
reproducibility record \citep{daskDelayed,daskArray}. Changing chunking or a
random-library version can change a pseudorandom realization even when the
statistical model and human-readable seed are unchanged.

The numerical tables and plotted coordinates in
\texttt{simulation-data/} are produced from the Dask-generated data.
XeLaTeX and PGFPlots render those values without Matplotlib. SciPy
supplies established distribution and numerical-integration routines;
the project does not implement a second normal-tail library.
The reproducibility record identifies the active Python source, model
configuration, seeds, and calculated assets.

Deal contracts describe selected input, output and side-effect
properties of the Python implementation. Runtime tests and statistical
diagnostics are not proofs of the pseudorandom generator or the
underlying population model. Theorems checked in Rocq apply only to
their explicitly stated mathematical hypotheses.

The main limitations are therefore specific. The normal generator is a
finite-state pseudorandom implementation; the normal model, fixed \(\rho\)
and stable reference population are assumptions; numerical integration has
separate accuracy checks; confidence intervals quantify finite sampling under
the model; and actual test validity, retesting behavior and real membership
proportions require evidence outside this simulation. The detailed counts,
replicate variation and calibrated-rule comparison show what the implemented
model predicts without eliminating those distinctions.

\appendix
\chapter{Referenced foundations and full working proofs}
\label{sec:proof-foundations}

The results below supply the mathematical foundations used in this report.
The explicitly assumed inputs are finite second moments and, where stated,
a jointly Gaussian law and independent observations. An assumption is not a
consequence of a simulation. Definitions, source quotations, historical
records, rounded numerical evaluations and recommendations have the status
specified in the claim map; they are not additional mathematical theorems.
The proofs reduce the deductions to linearity of integration, the spectral
theorem, the Gaussian characteristic function and conditional expectation.
Source locators refer to the linked author versions specified in the
bibliography, rather than silently mixing their pagination with print editions.

\noindent\textbf{Document-revision provenance.}
The current book-format revision is recorded in
\path{reporting/book-format/revision.json}, with current TeX/Bib sources
and PDF hashes. The earlier proof-text revision remains recorded in
\path{reporting/proof-revision/revision.json}.
The original \path{reporting/run.json} retains the numerical
simulation campaign, its original document hashes and numerical outputs.
The book revision checks and reuses the unchanged simulation-data files; it does
not rerun the experiment or claim a new native SPARK/Rocq proof.
The repository's subsequent shared-library refactor is separate from the
retained native campaigns. Their source-pinned results do not constitute
new GNATprove evidence for refactored Ada bodies; those changes require
their own matching campaign. The written mathematical proofs below do not
depend on that program-verification status.

\begin{proposition}[Finite moments and affine changes of scale]
\label{pf:moments}
Let a real random vector \(X\) have finite second moments, mean \(\mu\)
and covariance \(\Sigma\). Let \(Y\) be square-integrable, \(w\) a
deterministic vector of matching dimension, and \(b,s\) real constants. Then
\[
 \E(w^{\mathsf T}X)=w^{\mathsf T}\mu,\qquad
 \Var(w^{\mathsf T}X)=w^{\mathsf T}\Sigma w,\qquad
 \SD(b+sY)=|s|\SD(Y).
\]
A zero variance is equivalent to equality with the mean almost surely.
\end{proposition}
\begin{proof}
H\"older's inequality ensures integrability of all products
\((X_i-\mu_i)(X_j-\mu_j)\). Finite linearity of the integral then gives
the expectation formula and
\[
 \E\!\left[\Big(\sum_iw_i(X_i-\mu_i)\Big)^2\right]
 =\sum_{i,j}w_iw_j\E[(X_i-\mu_i)(X_j-\mu_j)]
 =w^{\mathsf T}\Sigma w.
\]
Replacing \(Y\) by \(b+sY\) multiplies its centered value by \(s\), so
its variance is \(s^2\Var(Y)\); taking the nonnegative square root proves
the SD formula. If \(\E[(Y-\E Y)^2]=0\), Markov's inequality applied to
\((Y-\E Y)^2\) makes \(P(|Y-\E Y|>\epsilon)=0\) for every
\(\epsilon>0\). The countable union over \(\epsilon=1/n\) proves the
almost-sure assertion. The reverse implication follows by integration.
These are applications of linearity, H\"older and the variance identities
in \citet[Section 1.6, especially Theorems 1.6.1, 1.6.3--1.6.4]{durrett2019}.
\end{proof}

\begin{proposition}[Covariance feasibility and Gaussian construction]
\label{pf:spectral}
A real symmetric matrix is a covariance matrix if and only if it is
positive semidefinite. For the equicorrelation matrix
\(R_n=(1-c)I_n+c\mathbf1_n\mathbf1_n^{\mathsf T}\), \(n\geq2\), this
is equivalent to \(-1/(n-1)\leq c\leq1\). Positive definiteness requires
both inequalities to be strict.
\end{proposition}
\begin{proof}
Proposition~\ref{pf:moments} gives
\(w^{\mathsf T}\Sigma w=\E[(w^{\mathsf T}(X-\mu))^2]\geq0\).
Conversely, the real spectral theorem diagonalizes a symmetric positive
semidefinite matrix as \(QDQ^{\mathsf T}\), with orthogonal \(Q\) and
nonnegative diagonal entries \(D\). For independent standard normal
coordinates \(Z\), the vector \(QD^{1/2}Z\) has covariance
\(QD^{1/2}I D^{1/2}Q^{\mathsf T}=\Sigma\) and is Gaussian, including
zero eigenvalues. The spectral theorem and positive square roots are
\citet[Theorems 7.29 and 7.38]{axler2024}; the Gaussian construction is
also \citet[Theorems 3.4.4--3.4.5]{gallager2013}.
For \(R_n\), the span of \(\mathbf1_n\) has eigenvalue
\(1+(n-1)c\), and its orthogonal complement has eigenvalue \(1-c\),
because \(\mathbf1_n^{\mathsf T}w=0\) there. Every vector decomposes
into these two orthogonal subspaces. Nonnegativity, or positivity, of
both eigenvalues gives precisely the stated inequalities. If one fixed
\(c\) must work for arbitrarily large \(n\), the lower bounds force
\(c\geq0\); the finite-dimensional bounds alone permit negative \(c\).
\end{proof}

\begin{proposition}[Gaussian closure and block independence]
\label{pf:gaussian}
An affine image of a jointly Gaussian vector is jointly Gaussian. Its law
is determined by its mean and covariance. Two blocks of a jointly Gaussian
vector are independent whenever their cross-covariance is zero, including
singular Gaussian vectors.
\end{proposition}
\begin{proof}
For a Gaussian vector with mean \(\mu\) and covariance \(\Sigma\), its
characteristic function is
\(\psi_X(t)=\exp(i t^{\mathsf T}\mu-t^{\mathsf T}\Sigma t/2)\).
Thus \(\psi_{AX+b}(t)=e^{it^{\mathsf T}b}\psi_X(A^{\mathsf T}t)\)
has the Gaussian form with mean \(A\mu+b\) and covariance
\(A\Sigma A^{\mathsf T}\). This also follows from the independent-normal
affine representation, without invertibility. If a two-block covariance
has zero off-diagonal blocks, the displayed characteristic function
factors as \(\psi_U(s)\psi_V(t)\). By Fubini this is the characteristic
function of the product of the two marginal measures. Uniqueness of
characteristic functions proves both independence and determination by
mean and covariance. The Gaussian formula and closure are
\citet[Theorems 3.3.1--3.3.2, equations (3.14) and (3.18)]{gallager2013};
uniqueness is \citet[Theorem 3.10.4]{durrett2019}.
\end{proof}

\begin{proposition}[Gaussian conditional distribution, including exact histories]
\label{pf:conditioning}
Suppose \((U,V)\) is jointly Gaussian with mean \((\mu_U,\mu_V)\),
covariance blocks \(A,B\) and \(C=\Cov(U,V)\), and \(B\) positive
definite. A version of the conditional law, for every \(v\), is Gaussian
with mean and covariance
\[
 \mu_U+CB^{-1}(v-\mu_V),\qquad A-CB^{-1}C^{\mathsf T}.
\]
The conditional-mean formula is the unique continuous version on the
whole conditioning space. The full joint covariance need not be invertible.
\end{proposition}
\begin{proof}
Set \(L=CB^{-1}\) and
\(W=U-\mu_U-L(V-\mu_V)\). Its cross-covariance with \(V\) is
\(C-LB=0\). Expanding its covariance gives
\(A-LC^{\mathsf T}-CL^{\mathsf T}+LBL^{\mathsf T}
=A-CB^{-1}C^{\mathsf T}\). Proposition~\ref{pf:gaussian} makes
\(W\) Gaussian and independent of \(V\). Consequently the kernel
\(K(v,H)=P(\mu_U+L(v-\mu_V)+W\in H)\) satisfies
\[
 P(U\in H,V\in D)=\int_D K(v,H)P_V(dv)
\]
by independence and Fubini, which is the defining conditional-law identity.
Its expectation is the displayed affine function. A positive-definite
Gaussian \(V\) has a strictly positive density everywhere. If two
continuous versions differed at any point, continuity would give a
neighborhood of positive probability on which they differ, contradicting
almost-sure uniqueness. Conditioning on an exact continuous history uses
this specified version, not division by the probability-zero event
\(\{V=v\}\). This residual proof is the Gaussian regression theorem
\citep[Theorems 3.5.1--3.5.2]{gallager2013}, proved here in the slightly
more general form allowing singular \((U,V)\); conditional versions are
defined in \citet[Section 4.1]{durrett2019}.
\end{proof}

\begin{proposition}[Normal marginal distributions do not imply joint normality]
\label{pf:normal-marginals}
Normal marginals and a specified Pearson correlation do not in general
determine Gaussian composite or joint-tail laws.
\end{proposition}
\begin{proof}
Let \(X\sim N(0,1)\) and independently let \(S\) equal \(1\) with
probability \((1+\rho)/2\) and \(-1\) otherwise, \(-1<\rho<1\).
Put \(Y=SX\). Symmetry gives \(Y\sim N(0,1)\), and independence gives
\(\E[XY]=\E[S]\E[X^2]=\rho\). Nevertheless
\(P(X+Y=0)=(1-\rho)/2>0\). A nondegenerate normal variable has a
density and no atoms; \(X+Y\) has variance \(2+2\rho>0\), so this
pair is not jointly Gaussian. For a concrete tail nonidentification at
\(\rho=0\), let \(t>0\) and \(p=P(X>t)\in(0,1/2)\). This pair
has \(P(X>t,Y>t)=p/2\), whereas an independent standard-normal pair,
with the same marginals and zero correlation, has intersection \(p^2\).
These differ since \(p<1/2\). The construction uses symmetry,
independence and the product-integration rule
\citep[Sections 1.7 and 2.1]{durrett2019}; it is a counterexample,
not an alternative assumption secretly used for the report's Gaussian tables.
\end{proof}

\begin{proposition}[Event arithmetic and conditional denominators]
\label{pf:events}
For events \(A,B\) with probabilities \(p_A,p_B\) and intersection
\(J\), their union has probability \(p_A+p_B-J\). If its conditioning
probability is positive, \(P(B\mid A)=J/p_A\). The correlation of the
nonconstant indicators is
\[
 \frac{J-p_Ap_B}{\sqrt{p_A(1-p_A)p_B(1-p_B)}}.
\]
\end{proposition}
\begin{proof}
The pointwise identities
\(1_{A\cup B}=1_A+1_B-1_{A\cap B}\) and
\(1_A1_B=1_{A\cap B}\) prove the union and covariance formulas by
linearity. Because \(1_A^2=1_A\),
\(\Var(1_A)=p_A-p_A^2\), and similarly for \(B\); divide their
covariance by the product of their positive SDs. The conditional formula
is precisely the event-conditional probability definition, requiring
\(p_A>0\). The four disjoint cells have probabilities
\(J,p_A-J,p_B-J,1-p_A-p_B+J\), obtained by the same identities.
These are consequences of finite additivity and conditional probability
\citep[Sections 1.1 and 4.1]{durrett2019}.
\end{proof}

\begin{proposition}[Tower property and population--conditional variance]
\label{pf:expectation}
For square-integrable \(Y\) and a sub-\(\sigma\)-algebra \(\mathcal G\),
\[
 \E\E[Y\mid\mathcal G]=\E Y,\qquad
 \Var(Y)=\E\Var(Y\mid\mathcal G)
              +\Var(\E[Y\mid\mathcal G]).
\]
\end{proposition}
\begin{proof}
The defining conditional-integral identity with the whole space gives
the first formula. Write \(M=\E[Y\mid\mathcal G]\). Conditional
Jensen \citep[Theorem 4.1.10]{durrett2019} ensures \(M\) is
square-integrable. In
\(Y-\E Y=(Y-M)+(M-\E Y)\), the cross term has expectation zero:
the pull-out and tower properties give
\(\E[(Y-M)(M-\E Y)]
=\E[(M-\E Y)\E(Y-M\mid\mathcal G)]=0\).
Expanding the square and integrating gives the variance identity. The
pull-out rule extends from bounded multipliers by truncation and
Cauchy--Schwarz. These conditional-expectation properties follow from
\citet[Section 4.1, especially Theorems 4.1.13--4.1.14]{durrett2019}.
In particular, fixing an observed starting score and varying it across
the population are different probability experiments, even when both
use the same marginal test SD.
\end{proof}

\chapter{Statistical reference laws and uncertainty calculations}
\label{sec:shared-statistical-methods}
\label{sec:simulation-statistical-methods}

These results concern ideal independent observations from the stated model.
They justify reference laws and interval calculations, rather than proving
that finite-state pseudorandom output is independent or that the model fits
real tests. The input theorems are Gaussian affine closure and the Gaussian
characteristic function \citep[Theorems 3.3.1--3.3.2]{gallager2013},
orthonormal-basis extension \citep[Theorem 6.36]{axler2024}, and linearity,
independence and elementary probability inequalities
\citep[Sections 1.6, 2.1 and 3.4]{durrett2019}. The derivations below do
not depend on a successful software test.

\begin{proposition}[Normal sample moments and the chi-square pivot]
\label{sm:normal-sample}
For an integer \(m\geq2\), independent \(Y_i\sim N(\mu,v)\), \(v>0\),
the sample mean is \(N(\mu,v/m)\). With
\(S^2=\sum_i(Y_i-\bar Y)^2/(m-1)\),
\((m-1)S^2/v\sim\chi^2_{m-1}\), independently of \(\bar Y\), and
\(\E S^2=v\).
\end{proposition}
\begin{proof}
Finite linearity and zero cross-covariances give the mean and variance;
Gaussian affine closure gives the law. Set \(Z_i=(Y_i-\mu)/\sqrt v\),
and extend the unit vector \(m^{-1/2}(1,\ldots,1)\) to an orthonormal
basis. Let \(Q\) have this vector as its first row. The transformed
Gaussian \(T=QZ\) has covariance \(QQ^{\mathsf T}=I\); its characteristic
function factors, so all its coordinates are independent standard normals.
Its first coordinate is \(\sqrt m(\bar Y-\mu)/\sqrt v\). Orthogonality
preserves squared norm, hence
\[
 \frac{1}{v}\sum_i(Y_i-\bar Y)^2
 =\sum_iZ_i^2-T_1^2=\sum_{j=2}^mT_j^2.
\]
The right side is chi-square by definition, independent of \(T_1\).
Its expectation is \(m-1\), proving unbiasedness of \(S^2\).
This argument supplies the normal and chi-square laws directly from the
referenced Gaussian and linear-algebra theorems.
\end{proof}

\begin{corollary}[SD intervals and fixed-parameter reference bands]
\label{sm:sd-interval}
Let \(q_u\) denote a \(\chi^2_{m-1}\) quantile and \(0<\alpha<1\).
An exact \(1-\alpha\) confidence interval for the unknown SD under this
normal model is
\[
 \left[\sqrt{\frac{(m-1)S^2}{q_{1-\alpha/2}}},
       \sqrt{\frac{(m-1)S^2}{q_{\alpha/2}}}\right].
\]
For a fixed theoretical \(v\), a reference interval for \(S^2\) is instead
\([vq_{\alpha/2}/(m-1),vq_{1-\alpha/2}/(m-1)]\). The reference
mean discrepancy in standard-error units is
\((\bar Y-\mu)/\sqrt{v/m}\sim N(0,1)\).
\end{corollary}
\begin{proof}
The continuous chi-square law gives
\(P(q_{\alpha/2}\leq(m-1)S^2/v\leq q_{1-\alpha/2})=1-\alpha\).
Solving these inequalities for \(v\), then taking the increasing positive
square root, gives the SD interval. Solving for \(S^2\), keeping \(v\)
fixed, gives the reference interval. The last assertion is the standardized
sample-mean law in Proposition~\ref{sm:normal-sample}. If \(v=0\), every
observation equals \(\mu\) almost surely, so this division-based pivot is
inapplicable; constant output is checked directly. SciPy evaluates the
quantiles \citep{scipystats}; the algebra, rather than the library call,
establishes the probability statement.
\end{proof}

\begin{proposition}[Bernoulli counts and sampling precision]
\label{sm:binomial}
For an integer \(m\geq1\) and \(m\) independent event indicators of common
probability \(p\in[0,1]\), the
count \(K\) is binomial, \(\E K=mp\), and
\(\Var(K/m)=p(1-p)/m\).
\end{proposition}
\begin{proof}
Independence makes the probability generating function of their sum
\((1-p+pt)^m\). The binomial theorem gives
\(P(K=j)=\binom mjp^j(1-p)^{m-j}\). Each indicator has mean \(p\)
and variance \(p(1-p)\), because its square equals itself. Summing means
and independent variances, then dividing by \(m^2\), proves the claims.
These are finite-product independence and linearity
\citep[Sections 1.6 and 2.1]{durrett2019}.
\end{proof}

\begin{proposition}[Wilson score endpoints and their coverage scope]
\label{sm:wilson}
For an integer \(m\geq1\), an observed count \(K\in\{0,\ldots,m\}\),
\(0<\alpha<1\), \(\hat p=K/m\) and
\(z=\Phi^{-1}(1-\alpha/2)\), inversion
of the normal score inequality gives endpoints
\[
 \frac{\hat p+z^2/(2m)\ \mathord\pm\
 z\sqrt{\hat p(1-\hat p)/m+z^2/(4m^2)}}{1+z^2/m}.
\]
The endpoint algebra is exact; nominal \(1-\alpha\) coverage uses a normal
score approximation and is not an exact finite-sample coverage theorem.
\end{proposition}
\begin{proof}
Squaring \(|\hat p-p|\leq z\sqrt{p(1-p)/m}\) gives
\[
 (m+z^2)p^2-(2m\hat p+z^2)p+m\hat p^2\leq0.
\]
The positive leading coefficient means precisely the interval between its
two real roots. Its discriminant is
\(4mz^2\hat p(1-\hat p)+z^4\). The quadratic formula, dividing
numerator and denominator by \(2m\), gives the stated endpoints. For
\(0<p<1\), the classical independent-identically-distributed central
limit theorem implies
\(\sqrt m(\hat p-p)/\sqrt{p(1-p)}\Rightarrow N(0,1)\), giving limiting
coverage at its continuity points, not equality for fixed \(m\).
This is the score construction of \citet{wilson1927}; the limit theorem is
\citet[Theorem 3.4.1]{durrett2019}. The implementation reuses
SciPy's binomial-result interval routine \citep{scipyBinomtest}.
\end{proof}

\begin{proposition}[Exact equal-tail binomial intervals]
\label{sm:clopper-pearson}
For an integer \(m\geq1\), a count \(j\in\{0,\ldots,m\}\), and
\(0<\alpha<1\), define \(L(0)=0\), \(U(m)=1\); for \(j>0\),
let \(L(j)\) solve \(P_{L(j)}(K\geq j)=\alpha/2\), and for \(j<m\),
let \(U(j)\) solve \(P_{U(j)}(K\leq j)=\alpha/2\). Then
\(P_p(L(K)\leq p\leq U(K))\geq1-\alpha\) for every \(p\in[0,1]\).
\end{proposition}
\begin{proof}
For \(0<j\leq m\), differentiating the finite binomial sum and
telescoping gives
\[
 \frac{d}{dp}P_p(K\geq j)
 =m\binom{m-1}{j-1}p^{j-1}(1-p)^{m-j}>0\quad(0<p<1).
\]
The tail runs continuously from zero to one, so \(L(j)\) exists uniquely;
the lower tail is strictly decreasing and gives \(U(j)\) similarly.
For \(0<j<m\), at \(p=L(j)\) the lower tail is at least
\(1-\alpha/2>\alpha/2\); its decrease implies \(U(j)>L(j)\).
The endpoint counts are ordered directly, so these really are intervals.
For a fixed \(p\), \(L(j)>p\) implies \(P_p(K\geq j)<\alpha/2\).
Such \(j\)'s form an upper terminal set, because upper tails decrease
with \(j\). If nonempty its probability is the tail at its smallest
member and is at most \(\alpha/2\). The analogous lower terminal set
\(U(j)<p\) has probability at most \(\alpha/2\). The union bound makes
total noncoverage at most \(\alpha\); endpoint parameters follow directly.
Discreteness explains why the inequality can be strict. This is the
test-inversion construction of \citet{clopper1934}, with its coverage
proof supplied here. Equivalently the endpoints are beta quantiles:
the derivative above is the \(\mathrm{Beta}(j,m-j+1)\) density and its
tail has value zero at \(p=0\), so integration gives the same cumulative
distribution; the complementary tail gives the upper endpoint.
\end{proof}

\begin{proposition}[Conditional frequencies with random denominators]
\label{sm:conditional-count}
For independent repetitions of a four-cell event experiment, condition on
\(N_A=\ell>0\) observations in event \(A\), where \(P(A)>0\).
The count also in \(B\) is then \(\mathrm{Binomial}(\ell,P(B\mid A))\).
Thus a binomial interval conditional on its observed denominator has its
stated exact or approximate coverage scope. A zero denominator gives no
defined conditional-frequency estimate.
\end{proposition}
\begin{proof}
Let \(q=P(A)\) and \(r=P(B\mid A)\). Grouping cells into
\(A\cap B,A\setminus B,A^c\), their probabilities are
\(qr,q(1-r),1-q\). Product independence and the multinomial counting
coefficient give
\[
 P(N_{AB}=j,N_A=\ell)
 =\binom m\ell\binom\ell j q^\ell(1-q)^{m-\ell}
           r^j(1-r)^{\ell-j}.
\]
Divide by \(P(N_A=\ell)=\binom m\ell q^\ell(1-q)^{m-\ell}\)
when positive. The remaining factor is the stated binomial mass. Cases
with \(q=1\) follow directly with \(\ell=m\). Conditional coverage for
each positive \(\ell\) also gives coverage averaged over positive
denominators by total probability. No convention assigning a number to
\(0/0\) supplies data when \(\ell=0\).
\end{proof}

\begin{proposition}[Pooling counts and moments]
\label{sm:pooling}
For a nonempty finite collection of groups of integer sizes \(m_j\geq2\),
means \(\bar Y_j\) and sample variances
\(S_j^2\), let \(M=\sum_jm_j\) and
\(\bar Y=M^{-1}\sum_jm_j\bar Y_j\). The pooled sample variance satisfies
\[
 (M-1)S^2=\sum_j(m_j-1)S_j^2+\sum_jm_j(\bar Y_j-\bar Y)^2.
\]
Independent equal-design Bernoulli groups pool to a binomial count of size
\(M\); independence is required for that probability law, not for the
displayed algebraic identity.
\end{proposition}
\begin{proof}
Within group \(j\), write
\(Y_{ji}-\bar Y=(Y_{ji}-\bar Y_j)+(\bar Y_j-\bar Y)\).
Squaring and summing, its cross term vanishes because
\(\sum_i(Y_{ji}-\bar Y_j)=0\). Summing groups gives the identity.
For counts, independence multiplies their generating functions to
\(\prod_j(1-p+pt)^{m_j}=(1-p+pt)^M\), proving the binomial law by
Proposition~\ref{sm:binomial}. Sharing random numbers across cases does not
establish this independence premise; in the overlapping-factor designs used
here the rows cannot be counted as additional independent replications.
Averaging group SDs generally does
not satisfy the displayed within-plus-between identity.
\end{proof}

\begin{proposition}[Simultaneous diagnostics and rare-event precision]
\label{sm:multiple}
If \(J\) fixed diagnostics have false-alarm probabilities at most
\(\alpha_j\), the probability of any false alarm is at most
\(\sum_{j=1}^J\alpha_j\), regardless of their dependence.
\end{proposition}
\begin{proof}
The pointwise bound \(1_{\cup_j F_j}\leq\sum_j1_{F_j}\), integrated,
is the union bound. Taking each \(\alpha_j=\alpha/J\) therefore gives
a family bound \(\alpha\). The same argument applies to two tails per
row, counting both in \(J\). Finite diagnostics remain observations
under a reference model and cannot establish an entire distribution.
The bound is the elementary subadditivity theorem
\citep[Section 1.1]{durrett2019}.
\end{proof}

\begin{corollary}[Zero counts, relative precision and sample-size scaling]
\label{sm:rare}
For independent Bernoulli trials, \(P(K=0)=(1-p)^m\). For \(p>0\) the
relative standard error of \(K/m\) is \(\sqrt{(1-p)/(mp)}\); requiring
it at most \(\epsilon>0\) means \(m\geq(1-p)/(p\epsilon^2)\).
At fixed \(p\in(0,1)\), multiplying \(m\) by four halves the standard
error. Zero successes do not prove \(p=0\).
\end{corollary}
\begin{proof}
The first statement is the binomial mass at zero. Divide the SD in
Proposition~\ref{sm:binomial} by \(p\), square the nonnegative inequality,
and rearrange to get the sample-size bound. The square-root dependence
proves the scaling. Every \(0<p<1\) assigns positive probability to
zero successes for finite \(m\), proving the final statement. A desired
95\% zero-count upper confidence bound, for example, follows from
\((1-U)^m=.025\) for the equal-tail construction, not from setting
\(p\) equal to zero. These conclusions concern sampling uncertainty;
they do not bound floating-point quadrature error or model error.
\end{proof}

% Project-specific proofs, included verbatim in each Towers report.
\section{Proofs for the multiple-test model and its numerical interpretation}
\label{sec:towers-proofs}

This appendix proves the deductive claims of the Towers study. Its probability
space is the ideal model, rather than a finite pseudorandom generator or an
empirical population. The foundational results and their proofs are in
Appendix~\ref{sec:proof-foundations}; the sampling results and their proofs are
in Appendix~\ref{sec:shared-statistical-methods}. The underlying standard
results are change of variables, Tonelli's theorem, dominated convergence,
the fundamental theorem of calculus and the intermediate value theorem
\citep{durrett2019,gallager2013}. Numerical tables are approximate evaluations
of the exact expressions below. Source dates, source quotations, software
inventories and retained execution counts are records to inspect, rather
than conclusions of probability theory.

\begin{proposition}[The Gaussian construction, density and joint tail]
\label{tw:gaussian}
Let \(G,E\) be independent standard normal variables, let \(|\rho|<1\), and
put \(X=G\), \(Y=\rho G+\sqrt{1-\rho^2}E\). Then \(X,Y\) have zero means,
unit variances, correlation \(\rho\), and joint density
\begin{equation}
 f_\rho(x,y)=\frac{1}{2\pi\sqrt{1-\rho^2}}
 \exp\!\left[-\frac{x^2-2\rho xy+y^2}{2(1-\rho^2)}\right].
 \label{eq:tw-density}
\end{equation}
For finite \(a,b\), their joint exceedance probability is
\begin{equation}
 J(a,b;\rho)=\int_a^\infty\varphi(x)
 \overline\Phi\!\left(\frac{b-\rho x}{\sqrt{1-\rho^2}}\right)\,dx.
 \label{eq:tw-tail-proof}
\end{equation}
For \(0<p<1\), the unique cutoff with exceedance probability \(p\) is
\(\Phi^{-1}(1-p)\). Strict and non-strict versions have equal probability.
\end{proposition}
\begin{proof}
Linearity gives the two means. Independence gives \(E[GE]=E[G]E[E]=0\);
expansion of the squares gives
\(\operatorname{Var}(Y)=\rho^2+(1-\rho^2)=1\) and
\(\operatorname{Cov}(X,Y)=\rho\). Gaussian affine closure is proved in
Proposition~\ref{pf:gaussian}. Alternatively the invertible transformation has
inverse \(g=x\), \(e=(y-\rho x)/\sqrt{1-\rho^2}\) and Jacobian determinant
\((1-\rho^2)^{-1/2}\). Multiplying \(\varphi(g)\varphi(e)\) by that
determinant and expanding the exponent proves \eqref{eq:tw-density}.
Dividing it by the positive marginal density \(\varphi(x)\) gives the
conditional density of \(Y\) at fixed \(X=x\): it is normal with mean
\(\rho x\) and variance \(1-\rho^2\). This is the continuous conditional
version explained in Proposition~\ref{pf:conditioning}; an event \(X=x\) itself
has probability zero. Tonelli's theorem integrates its upper tail against
\(\varphi(x)\), proving \eqref{eq:tw-tail-proof}. Finally \(\Phi\) is
continuous, strictly increasing, and has limits zero and one, because
\(\varphi>0\) and integrates to one. Inverting it proves the cutoff claim.
Its continuous density assigns each cutoff point probability zero.
\end{proof}

\begin{proposition}[Every denominator and the binary correlation]
\label{tw:events}
Write \(A=\{X>a\}\), \(B=\{Y>b\}\), \(p_A=P(A)\), \(p_B=P(B)\),
and \(J=P(A\cap B)\). The four disjoint cell probabilities, in the order
neither, \(A\) only, \(B\) only, both, are
\begin{equation}
 1-p_A-p_B+J,\quad p_A-J,\quad p_B-J,\quad J.
 \label{eq:tw-cells}
\end{equation}
Thus \(u=P(A\cup B)=p_A+p_B-J\), and
\begin{align}
 P(B\mid A)&=J/p_A,&P(A\mid B)&=J/p_B,\\
 P(A^c\mid B)&=1-J/p_B,&
 P(A\cap B\mid A\cup B)&=J/u.
 \label{eq:tw-denominators}
\end{align}
Here all denominators are positive for finite Gaussian cutoffs.
The Pearson correlation of the indicators is
\begin{equation}
 \operatorname{Corr}(1_A,1_B)=
 \frac{J-p_Ap_B}{\sqrt{p_A(1-p_A)p_B(1-p_B)}}.
 \label{eq:tw-phi-proof}
\end{equation}
At equal margins \(p\), exactly one pass has probability \(2(p-J)\),
one specified pass without the other has probability \(p-J\), and the
both-pass fraction of the eligible union is \(J/(2p-J)\).
\end{proposition}
\begin{proof}
The four cells partition the whole space. Subtracting the both-pass cell
from each marginal and then subtracting the other three cells from one
proves \eqref{eq:tw-cells}. Their nonnegativity proves the Fr\'echet bounds
\(\max(0,p_A+p_B-1)\leq J\leq\min(p_A,p_B)\), and summing the three
passing cells proves the union formula. In every displayed conditional,
intersect the numerator event with the conditioning event and divide by
the probability of that conditioning event, as in Proposition~\ref{pf:events}.
Since \(1_A^2=1_A\), its mean and variance are \(p_A\) and
\(p_A(1-p_A)\); likewise for \(B\). Also \(1_A1_B=1_{A\cap B}\), so
the covariance is \(J-p_Ap_B\). The correlation definition now proves
\eqref{eq:tw-phi-proof}. Substitution of equal margins proves the last
three assertions. These calculations do not identify binary correlation
with \(\rho\); each is a correlation of different random variables.
\end{proof}

\begin{proposition}[Correlation differentiation and sensitivity]
\label{tw:plackett}
For finite cutoffs and \(-1<\rho<1\),
\begin{equation}
 \frac{\partial J(a,b;\rho)}{\partial\rho}=f_\rho(a,b)>0,
 \qquad
 J(a,b;\rho)=\overline\Phi(a)\overline\Phi(b)
                 +\int_0^\rho f_t(a,b)\,dt.
 \label{eq:tw-plackett-proof}
\end{equation}
Consequently joint passing and either conditional passing probability
strictly increase with \(\rho\) at fixed marginal cutoffs, whereas the
union probability strictly decreases. Equal cutoffs give equal theoretical
conditionals; mixed cutoffs need not.
\end{proposition}
\begin{proof}
This is the bivariate case of Plackett's identity \citep{plackett1954}; a
derivation is included. With \(D=1-t^2\), direct differentiation gives
\[
 \partial_t f_t(x,y)=
 \left\{\frac{t}{D}+\frac{(1+t^2)xy-t(x^2+y^2)}{D^2}\right\}f_t(x,y)
 =\partial_x\partial_y f_t(x,y).
\]
On every compact correlation interval \(|t|\leq r<1\), the derivative is
bounded by a constant times
\((1+x^2+y^2)\exp[-(x^2+y^2)/(2(1+r))]\), an integrable function.
Dominated convergence therefore permits differentiation of the double
tail integral. Integrating \(\partial_x\partial_y f_t\) over
\([a,\infty)\times[b,\infty)\) gives \(f_t(a,b)\): the Gaussian
density and its first derivatives vanish at the infinite boundaries.
At \(t=0\) the density factors, so the joint tail is the product of the
marginal tails. The fundamental theorem of calculus gives the second
identity, including negative \(\rho\) with its oriented integral.
The density is strictly positive. The monotonicity claims follow by
dividing by the fixed positive margins, or subtracting \(J\) from their
sum. Swapping \(x,y\) leaves \(f_\rho\) unchanged, which proves the
equal-cutoff symmetry. At mixed margins the conditional ratio is exactly
\((J/p_A)/(J/p_B)=p_B/p_A\), explaining the factor of twenty in the
top-2\% versus top-0.1\% example.
\end{proof}

\begin{proposition}[A distribution-free contradiction and a Gaussian bound]
\label{tw:original-bound}
Marginal pass rates \(0.02\) and \(0.001\) cannot have union probability
\(0.0211\). Moreover, equal top-2\% Gaussian cutoffs at \(\rho=0.7\)
satisfy
\begin{equation}
 J>\frac{47}{15000}>0.0031,
 \qquad P(B\mid A)>\frac{47}{300}>0.156,
 \qquad P(A\cup B)<\frac{553}{15000}<0.0369.
 \label{eq:tw-rigorous-rejection}
\end{equation}
These exact inequalities already reject the first letter's overlap
\(0.002\), repeat-pass \(0.1027\) and union \(0.0383\) as model values,
without relying on the last digits of a numerical integrator.
\end{proposition}
\begin{proof}
The union bound gives \(P(A\cup B)\leq0.021<0.0211\).
For the Gaussian claims write \(T(x)=\overline\Phi(x)\). Integration of
\(t\varphi(t)=-\varphi'(t)\) gives the strict Mills upper bound
\(T(x)<\varphi(x)/x\) for \(x>0\), since \(t/x>1\) for \(t>x\).
For completeness the constant bound \(\pi>25/8\) also has an elementary
proof. The exact finite geometric identity
\[
 \frac1{1+t^2}=1-t^2+t^4-t^6+\frac{t^8}{1+t^2}
\]
and integration on \([0,x]\), \(x>0\), give
\(\arctan x>x-x^3/3+x^5/5-x^7/7\), because the remainder integral is
positive. The tangent addition formula gives
\(\arctan(1/2)+\arctan(1/3)=\pi/4\): its tangent is one and the
sum is in \((0,\pi/2)\), since its cosine has the positive factor
\(1-(1/2)(1/3)\). Applying the lower bounds at \(1/2\) and \(1/3\)
therefore gives
\[
 \frac\pi4>
 \sum_{x\in\{1/2,1/3\}}\left(x-\frac{x^3}{3}+\frac{x^5}{5}-\frac{x^7}{7}\right)
 =\frac{25}{32}+\frac{7765}{1959552}>\frac{25}{32}.
\]
Together with
\(e^{8/3}>\sum_{n=0}^4(8/3)^n/n!=3035/243>11\), this implies
\[
 T(7/3)<\frac{\varphi(7/3)}{7/3}
 <\frac{(2/5)}{11(7/3)}=\frac6{385}<\frac1{50}.
\]
Thus the cutoff \(a\) with \(T(a)=1/50\) is below \(7/3\).
Taylor's theorem with its signed third derivative gives
\(e^{-v}\leq1-v+v^2/2\) for \(v\geq0\). Hence
\[
 T(1)=\frac12-\int_0^1\varphi(t)\,dt
 >\frac12-\frac25\left(1-\frac16+\frac1{40}\right)
 =\frac{47}{300}.
\]
For every \(x\geq a\), the argument of the conditional tail in
\eqref{eq:tw-tail-proof} is at most
\(0.3a/\sqrt{0.51}<0.7/\sqrt{0.51}<1\). Therefore
\(J>T(a)T(1)>47/15000\). Division by \(1/50\) and the equal-margin
union identity prove the remaining inequalities.
\end{proof}

\begin{proposition}[Same-score nesting and different-test non-nesting]
\label{tw:nesting}
If \(b>a\) are applied to the same score \(Z\), then
\(\{Z>b\}\subset\{Z>a\}\). When \(P(Z>b)>0\), the conditional
probability of passing the lower cutoff given the higher is one.
With different Gaussian scores
and \(|\rho|<1\), neither finite-cutoff pass event contains the other
up to null sets.
\end{proposition}
\begin{proof}
The first inclusion is the elementary implication \(z>b\Rightarrow z>a\).
Its conditional probability is therefore the numerator probability divided
by itself. For different scores the open rectangles
\((a,a+1)\times(b-2,b-1)\) and
\((a-2,a-1)\times(b,b+1)\) lie respectively in \(A\setminus B\) and
\(B\setminus A\). The density \eqref{eq:tw-density} is continuous and
positive on their closures, so each rectangle has positive probability.
This proves non-nesting, strict union enlargement over either margin, and
strictly subunit conditional probabilities. The statement uses finite
cutoffs and a nonsingular Gaussian distribution; it does not assert a
relationship between actual admission policies or actual tests.
\end{proof}

\begin{proposition}[Existence and uniqueness of OR calibration]
\label{tw:calibration}
For every \(-1<\rho<1\) and \(0<u_0<1\), there is exactly one finite
\(t\) satisfying
\(2\overline\Phi(t)-J(t,t;\rho)=u_0\).
On a reporting scale \(I=100+16X\), its cutoff is \(100+16t\).
\end{proposition}
\begin{proof}
Let \(U(t)=P(X>t\text{ or }Y>t)\). Inclusion-exclusion identifies it with
the left side. If \(t_1<t_2\), the event at \(t_2\) is contained in the
event at \(t_1\). A rectangle with \(t_1<X<t_2\) and \(Y<t_1\) has
positive probability and belongs to the set difference; consequently
\(U(t_1)>U(t_2)\). For \(t_n\to t\), the event indicators converge
except on \(X=t\) or \(Y=t\), which have probability zero. Dominated
convergence proves continuity. As \(t\to\infty\), the union bound gives
\(U(t)\leq2\overline\Phi(t)\to0\). As \(t\to-\infty\),
\(U(t)\geq\overline\Phi(t)\to1\). The intermediate value theorem gives
existence and strict decrease gives uniqueness. A positive affine reporting
scale transforms the threshold by the same affine formula; its mean and
standard deviation are 100 and 16 by Proposition~\ref{pf:moments}.
\end{proof}

\begin{proposition}[A rarity equivalent identifies a size, not a set]
\label{tw:rarity}
For \(0<p<1\), \(100+16\Phi^{-1}(1-p)\) is the unique SD-16 normal
cutoff exceeded with probability \(p\). Equal selected fractions do not
imply equal selected people or identify any latent ability.
Normal margins and their correlation do not identify the Gaussian overlap.
\end{proposition}
\begin{proof}
The cutoff claim follows by inversion in Proposition~\ref{tw:gaussian}
and the positive affine conversion. To disprove set identification, let
\(U\) be uniform on \([0,1)\). The intervals \([0,p)\) and the same
interval translated modulo one by \(d=\min(p,1-p)/2>0\) both have
probability \(p\), but their symmetric difference has positive length.
Thus equal probability does not identify the event. For latent-variable
identification, keep the same observed pair \((X,Y)\) and the same observed
selection event \(A=\{X>a\}\). In one model
adjoin an independent standard normal variable \(L\), called latent ability.
In another model set \(L=X\). Both have the identical observed score law
and a standard normal \(L\), but observed selection is independent of
\(L\) in the first model and associated with it in the second. No
measurement relation specifying \(L\) was part of the observed model, so
the observations cannot choose between them. Finally
Proposition~\ref{pf:normal-marginals} constructs normal margins with the same
correlation using \(Y=SX\) with an independent biased sign \(S\).
For positive cutoffs its joint tail is
\((1+\rho)\overline\Phi(\max(a,b))/2\), not generally
\eqref{eq:tw-tail-proof}. This provides the stated explicit failure of
identification by margins and correlation alone.
\end{proof}

\begin{proposition}[Repeated routes and selection]
\label{tw:attempts}
For any finite pass events \(A_1,\ldots,A_m\),
\[
 \max_iP(A_i)\leq P(\cup_iA_i)\leq\min(1,\sum_iP(A_i)).
\]
If they are independent with common probability \(p\), the union is
\(1-(1-p)^m\). A population correlation does not, by itself, determine a
correlation in a selected subpopulation.
\end{proposition}
\begin{proof}
Each individual event is contained in the union, proving the lower bound.
Pointwise \(1_{\cup_iA_i}\leq\sum_i1_{A_i}\); expectations and the
universal bound one prove the upper bound. Independence of the complements
gives \(P(\cap_iA_i^c)=(1-p)^m\), and taking the complement proves the
formula. For a concrete selection counterexample take independent standard
normals \(X,Y\), and condition on \(XY>0\). This event has probability
one half; symmetry gives zero conditional means. Conditional squares retain
mean one, but the conditional product has expectation
\(E[|X|]E[|Y|]=2/\pi>0\). Indeed the signs are independent of the
absolute values, and \(E|X|=2\int_0^\infty x\varphi(x)\,dx
=\sqrt{2/\pi}\). The selected correlation is consequently \(2/\pi\),
despite population correlation zero. This is an existence counterexample,
not a model of any named society's selection.
\end{proof}

\begin{proposition}[What the BASIC grid normalizes]
\label{tw:grid}
Let \(h=1/2\), let \(F=\{1,1+h,\ldots,199\}^2\), and give node
\((x,y)\) weight \(g(x,y)>0\) proportional to its Gaussian IQ-scale
density. The ratio of the sum over passing nodes to the sum over \(F\)
is exactly a probability for a finite weighted distribution on \(F\).
It is not exactly a continuous Gaussian tail. There are \(397^2=157609\)
nodes; constant density-scale factors and \(h^2\) cancel from the ratio.
Viewed as composite midpoint quadrature, nodes starting at \(Q\) cover
cells starting at \(Q-h/2=Q-0.25\), not at \(Q\).
\end{proposition}
\begin{proof}
The normalized node masses are
\(w_{xy}=g(x,y)/\sum_Fg\); they are positive and sum to one. Summing
them over the passing subset proves the discrete-probability assertion and
also supplies its own consistent marginal denominator. The node count in
each coordinate is \((199-1)/(1/2)+1=397\). Multiplying every weight by
any common positive constant leaves every normalized mass unchanged,
proving cancellation of the cell area and the density Jacobian factor
\(16^{-2}\). Centering a cell of width \(h\) on its node gives the
stated lower boundary and upper boundary \(199+h/2\).

More precisely, let \(N,D\) be the continuous Gaussian masses of the
passing cell rectangle and the full cell rectangle, and let \(n,d\) be
their midpoint sums including cell area. On a finite rectangle a twice
continuously differentiable density satisfies the bound
\[
 |n-N|\leq\frac{|R|h^2}{24}(M_{xx}+M_{yy}),
\]
where \(M_{xx},M_{yy}\) bound the absolute second derivatives on the
rectangle. Apply the midpoint remainder first in \(x\), then in \(y\):
the linear Taylor term integrates to zero and
\(\int_{-h/2}^{h/2}u^2/2\,du=h^3/24\); summing the cells proves the
bound. The same argument bounds \(|d-D|\). Since \(d>0\),
\[
 \left|\frac nd-\frac ND\right|
 \leq\frac{|n-N|+(N/D)|d-D|}{d}.
\]
The truncated continuous ratio \(N/D\) is itself a conditional probability
given the full finite cell rectangle. Its excluded mass is bounded by the
sum of the four marginal tails beyond its edges, by the union bound.
This separates discretization, shifted boundary and truncation; it does
not claim that their errors vanish or have the same sign at finite \(h\).
Comparisons of their sizes in the tables are numerical comparisons, not a
certified integration-error bound. An integer cutoff \(Q\) has actual
continuous marginal \(\overline\Phi((Q-100)/16)\); division by a
different nominal tail answers no conditional probability of this model.
Rounding the scores first would define yet another event and requires an
explicit rounding convention.
\end{proof}

\begin{proposition}[Quantization of a conditional ratio]
\label{tw:quantization}
If \(|j'-j|\leq\delta\), \(|p'-p|\leq\delta\), \(p>\delta\),
and \(j\geq0\), then
\[
 \left|\frac{j'}{p'}-\frac jp\right|
 \leq\frac{\delta(1+j/p)}{p-\delta}.
\]
If also \(0\leq j\leq p\), this is at most
\(2\delta/(p-\delta)\).
\end{proposition}
\begin{proof}
Put \(\Delta j=j'-j\), \(\Delta p=p'-p\). The difference is
\((p\Delta j-j\Delta p)/(pp')\), and \(p'\geq p-\delta>0\).
The triangle inequality bounds the numerator by \(p\delta+j\delta\).
Division gives the first bound, and \(j/p\leq1\) gives the second.
For \(\delta=0.51\times10^{-15}\) and \(p=1/33000\), substituting
in the latter expression gives less than \(3.37\times10^{-11}\).
This accounts for probability quantization only, not integration error.
\end{proof}

\begin{proposition}[Simulation partition and selected denominators]
\label{tw:simulation}
For \(T\) simulated pairs, let \(n_{rs}\) count the four indicator
values \((1_A,1_B)=(r,s)\). Then
\[
 \sum_{r,s}n_{rs}=T,\quad n_A=n_{10}+n_{11},\quad
 n_B=n_{01}+n_{11},\quad n_U=n_A+n_B-n_{11}.
\]
Under ideal independent identical draws, each fixed event count is binomial.
Given \(n_A=m>0\), \(n_{11}\) is binomial with \(m\) trials and
success probability \(J/p_A\). Different cutoff rules evaluated on the
same pairs need not give independent estimates.
\end{proposition}
\begin{proof}
For each pair exactly one of the four Boolean outcomes occurs. Summing that
pointwise partition proves the integer identities without any distributional
assumption. Under the additional independent-pair model, a fixed event gives
independent Bernoulli indicators, whose sum has the binomial law by
Proposition~\ref{sm:binomial}. Conditional on the identities of the \(m\)
observations passing \(A\), their \(B\) indicators are independent with
conditional probability \(J/p_A\). Every choice of these identities gives
the same binomial sum; mixing over the choices proves the claim given
\(n_A=m\), as proved generally in Proposition~\ref{sm:conditional-count}.
When \(m=0\), the conditional frequency is undefined. To see dependence
between common-pair rules explicitly, if one pass event \(C\) is contained
in another \(D\) and both have probabilities strictly between zero and
one, their one-pair indicators have covariance
\(P(C)(1-P(D))>0\). Their frequency estimates from \(T\) independent
pairs have this covariance divided by \(T\), by expanding the double
sum and canceling distinct-pair covariances. Thus independence between rows
cannot be inferred from independent observations within each row.
\end{proof}

\chapter{Claim-to-proof and simulation-evidence register}
\label{sec:towers-simulation-claim-map}
Every numerical display is a rounded evaluation or a retained execution
record. The probability results concern independent draws from the ideal
model; the pseudorandom implementation is evaluated separately. This register
maps all sections and their mathematical assertions to the full proofs in
this report, and marks the assertions that instead require provenance or
empirical evidence.

\small
\begin{longtable}{>{\raggedright\arraybackslash}p{.29\linewidth}>{\raggedright\arraybackslash}p{.64\linewidth}}
\toprule
Claim or output & Proof or evidence status\\
\midrule\endhead
Experiment, model input and five rules
& Zero means, unit SDs, \(\rho=0.7\), stable margins and independent ideal
pairs are model inputs. Proposition~\ref{tw:gaussian} derives the generator
law, quantile thresholds and equality of strict/non-strict probabilities.
Seeds, chunking, version pins and trial counts are execution specifications,
not a proof that deterministic streams are independent random variables.\\
Union \eqref{eq:towers-union} and conditional ratios
\eqref{eq:towers-conditionals}
& Proposition~\ref{tw:events} derives every numerator and denominator by
the four-cell partition and the definition of conditional probability.\\
Generator \eqref{eq:towers-generator}; mean, variance, covariance
& Proposition~\ref{tw:gaussian} and Propositions~\ref{pf:moments},
\ref{pf:gaussian} prove the affine Gaussian law and moments. These are
population quantities, not assertions that each finite sample has those
exact moments.\\
Four cell counts and event masks
& Proposition~\ref{tw:simulation} proves the integer identities observation
by observation. The numerical occurrence of each retained count is an
execution record; count consistency does not prove the model.\\
Reference integral \eqref{eq:towers-reference-integral} and correlation integral
& Propositions~\ref{tw:gaussian}, \ref{tw:plackett} prove both exact
integral identities, including differentiation justification. SciPy computes
approximations of these targets. Their cross-check is numerical validation.\\
Pooled event table and Figure~\ref{fig:towers-events}
& Events have ideal binomial counts by Proposition~\ref{sm:binomial}.
Propositions~\ref{sm:pooling}, \ref{sm:wilson}, \ref{sm:clopper-pearson}
prove pooled estimation and interval construction. Frequencies and plot
coordinates come from the recorded integer counts, not an analytic theorem.
The exact E2/E0.1/Rare/Mixed targets are the preceding integrals at the
listed quantiles; printed decimals approximate them.\\
Enlarged eligibility and mixed non-nesting
& Proposition~\ref{tw:nesting} proves positive probability of each
specified-only outcome, hence strictly enlarged unions. Proposition~\ref{tw:events}
proves the mixed conditional factor twenty and complement rates.\\
Conditional table; fluctuating denominators and equal-cutoff symmetry
& Propositions~\ref{tw:events}, \ref{tw:plackett}, \ref{tw:simulation}
and Proposition~\ref{sm:conditional-count} prove symmetry of the targets and
the binomial law conditional on the actual selected denominator. A zero
denominator leaves the estimate undefined.\\
Rare case, relative precision, zero counts and replicate variation
& Proposition~\ref{sm:binomial} and Corollary~\ref{sm:rare} derive \(E[K]=Tj\),
\(\operatorname{Var}(K)=Tj(1-j)\), relative SE
\(\sqrt{(1-j)/(Tj)}\), and \(P(K=0)=(1-j)^T\).
Expected counts are exact substitutions of the integral; recorded rare
counts, extrema and standardized discrepancies are finite observations.\\
Pointwise versus simultaneous intervals; common random numbers
& Proposition~\ref{sm:multiple} proves the union-bound distinction.
Proposition~\ref{tw:simulation} proves possible nonzero covariance between
same-pair estimates. No independent-row standard error is inferred. Wilson
coverage is approximate; Proposition~\ref{sm:clopper-pearson} proves conservative
binomial coverage, not protection against model misspecification.\\
OR2 \eqref{eq:towers-calibration}; threshold, scale and achieved rate
& Proposition~\ref{tw:calibration} proves continuous strict monotonicity,
endpoint limits and the unique root. The generated calibration table is
the fixed-root simulation record. A 95\% interval can miss its target;
Proposition~\ref{sm:multiple} separates pointwise intervals from broad diagnostic
checks. The observed miss is retained without seed or cutoff tuning.\\
SPARK generator and ideal normal transformation
& \Cref{tw:box-muller} proves the continuous-uniform Box--Muller law by the
change-of-variables theorem and independence. \Cref{tw:gaussian} proves the
correlated pair's ideal law. The Park--Miller recurrence and finite checked
bounds are source-pinned implementation facts, not proofs of IID randomness.\\
All native cells, rates, campaign replay and Dask comparison
& \Cref{sec:towers-spark-study} identifies the immutable 31 October and
byte-identical 4 November outputs, source manifest and CSV hash. Integer
cells and execution statuses are recorded observations. \Cref{tw:simulation}
proves the cell-derived rates; \cref{sm:binomial,sm:conditional-count,sm:rare}
prove their ideal sampling laws. Native and Dask denominators are retained
separately. The fourfold-size/half-SE statement follows by substituting
\(4T\) into \(p(1-p)/T\), not from an assumption of independent finite
streams. Wilson and exact intervals have the distinct proofs in
\cref{sm:wilson,sm:clopper-pearson}.\\
Native proof totals, trusted completions and source context
& The dated 4 November portfolio and runtime manifests record 696 checks,
230 flow/466 prover, 99 project subprograms/557 clauses, exact replay and
\(10^{-10}\) numerical tolerance. These are source-pinned evidence,
not newly rerun proofs of later refactored source, Dask, third-party
libraries, ideal IID assumptions, physical I/O or model adequacy.\\
Binary correlation \eqref{eq:towers-phi} and phi table
& Proposition~\ref{tw:events} proves the population indicator formula from
its moments. Replacing probabilities by empirical frequencies gives the
sample indicator correlation when both sample variances are positive.
The table's empirical values are recorded arithmetic evaluations.\\
Historical claims, interpretations, latent ability and real membership
& \citet{towersOriginal,towersFollowup} supply source attributions.
Propositions~\ref{tw:original-bound}, \ref{tw:rarity}, \ref{tw:attempts}
give mathematical contradictions or explicit nonidentifiability
counterexamples. The model supplies no empirical membership or validity
study; recommendations remain recommendations.\\
Dask, libraries, contracts, SPARK, provenance and AI disclosure
& Source, versioned environment, generated manifest and dated native
evidence provide operational records. These are distinct from the written
model proofs, structural contract inventories, runtime/property tests,
limited symbolic proofs and trusted-completion boundaries. The report adds
no native SPARK campaign. No theorem here proves finite generator quality,
library accuracy, file persistence, actual test validity or human review.\\
\bottomrule
\end{longtable}
\normalsize


\bibliographystyle{plainnat}
\bibliography{references,simulation_references}
\chapter{AI disclosure and provenance}\label{sec:ai-disclosure}
% Compatibility alias: canonical AI-use disclosure is in library/latex.
% Canonical full appendix text; the chapter heading is defined in report_style.tex.
This report and its supporting repository were substantially generated and
revised with OpenAI Codex and collaborating AI agents. The session identifies
Codex as based on the GPT-6 family. Exact serving checkpoints, per-agent model
identifiers and generation settings were not preserved; no more precise model
attribution is claimed. This disclosure covers the mathematical exposition,
simulations, programs, documentation, verification, typesetting and publication,
including subsequent layout and documentation revisions.

AI contributions included interpreting the problem and historical articles;
developing and reviewing derivations, counterexamples, theorem references,
proofs and claim maps; calculating examples and limits; editing
Python/Dask/Deal implementations, contracts, tests and auditors; executing
the recorded checks; producing tables and scientific charts; preparing the
XeLaTeX/KOMA-Script books; visually reviewing PDFs; and preparing repository
updates.

The requester Anders Hellström supplied questions and sources, clarified notation and
averaging, selected example parameters, specified formats and verification
requirements, requested revisions and authorized publication. Repository
ownership and commit attribution do not identify every passage as human-authored
or independently verified. Parallel mathematical and visual reviews were
performed by AI agents and automated tools within the same workflow.
Independent human peer review, professional sign-off and endorsement by
original authors, source hosts or cited organizations are not recorded and
are not claimed.

Original authors and references remain credited. Historical articles are
analysed rather than relicensed; replacement articles are AI-generated
analytical reconstructions. Sources were consulted at different levels,
including mathematical texts, articles, official software documentation and
bibliographic metadata. Reference notes identify versions and access limits;
full-text examination of every cited work is not claimed. Mathematical proofs
establish consequences of stated hypotheses, rather than empirical adequacy
of the IQ model.

Verification claims rest on retained commands, actual outputs, tool versions,
source hashes and checksums. Previous disconnected implementations and
their proof campaigns were removed from maintained source; they remain accessible
through earlier Git history and are not current-source verification for
Python/Dask. Deal contract declarations, static lint, runtime/property
tests, exact Rocq kernel checks, numerical accuracy, PRNG properties and
model validity have separate scopes. Unsupported, skipped or unrun checks
are not reported as successful proofs. Trusted numerical/library boundaries
and specification adequacy remain explicit.
Independent numerical validation means a separate implementation, algorithm or
library, not independent human authorship or review.

No participant data were collected and no empirical psychometric validation
study was conducted.
Simulations use synthetic observations under stated probability models. A typesetting-only revision does not, by itself, execute additional random
simulations or prove additional Rocq theorems. Older
compiler evidence is available through Git history, not active source files.
The repository does not preserve a complete prompt/tool transcript, exact per-line authorship
or AI-contribution percentage; these have not been invented. Its license
applies to original material, subject to third-party rights.


The Python/Dask-focused editorial revision removes the discontinued simulation implementation and its comparisons; it does not assert any newly executed Monte Carlo campaign.
\end{document}
